% Extracto íntegro por residencia del cierre físico-matemático sucesor.
\section{Cristal temporel fini HMT : forçage, sous-harmonique et stabilité spectrale}
\label{sec:hmtf2-cristal-temporal}

\subsection{Forçage construit à partir de l’horloge à 108 états}

Soit \(N=108\),
\[
 \mathcal H_{\rm clk}=\mathbb C^N,\qquad
 S|j\rangle=|j+1\rangle,\qquad
 Z|j\rangle=\zeta^j|j\rangle,\qquad
 \zeta=\exp(2\pi{\rm i}/N).
\]
Dans la base de Fourier \(\{|\widetilde k\rangle\}_{k=0}^{N-1}\), nous fixons
\[
 H_{\rm clk}
 =\hbar_{\rm int}\frac{2\pi}{Nt_0}
 \sum_{k=0}^{N-1}k|\widetilde k\rangle\langle\widetilde k|,
 \qquad
 \exp\!\left(-\frac{{\rm i}}{\hbar_{\rm int}}H_{\rm clk}t_0\right)=S.
\]
Cette égalité fixe l’horloge de départ. Pour exhiber une impulsion périodique non
constante, soit \(0<\eta<1\), \(T_{\rm d}=t_0\), et
définissons, à chaque période,
\[
 H_{\rm d}(t)=
 \begin{cases}
  0,&0\leq t<\eta T_{\rm d},\\[2mm]
  (1-\eta)^{-1}H_{\rm clk},
  &\eta T_{\rm d}\leq t<T_{\rm d}.
 \end{cases}
 \qquad H_{\rm d}(t+T_{\rm d})=H_{\rm d}(t).
\]
Comme les deux segments commutent, leur propagateur de Floquet est
\[
 U_F
 =\mathcal T\exp\!\left(
 -\frac{{\rm i}}{\hbar_{\rm int}}\int_0^{T_{\rm d}}H_{\rm d}(t)\,\mathrm dt
 \right)
 =\exp\!\left(-\frac{{\rm i}}{\hbar_{\rm int}}H_{\rm clk}T_{\rm d}\right)
 =S.
\]

Nous définissons les observables hermitiens
\[
 Q=\frac{Z+Z^\dagger}{2},
 \qquad
 P=\frac{Z-Z^\dagger}{2{\rm i}}.
\]

\begin{theorem}[Sous-harmonique primitive du cristal temporel fini HMT]
\label{thm:hmtf2-subarmonia-108}
Pour \(\rho_{j_0}=|j_0\rangle\langle j_0|\), le vecteur d’ordre
\[
 \mathbf q_n
 =\left(
 \operatorname{Tr}\rho_{j_0}U_F^{-n}QU_F^n,\,
 \operatorname{Tr}\rho_{j_0}U_F^{-n}PU_F^n
 \right)
\]
satisfait
\[
 \mathbf q_n
 =\left(
 \cos\frac{2\pi(j_0+n)}N,\,
 \sin\frac{2\pi(j_0+n)}N
 \right).
\]
Sa période primitive est \(N T_{\rm d}=108T_{\rm d}\), tandis que le forçage
a pour période \(T_{\rm d}\). Le système réalise donc une réponse
sous-harmonique primitive d’ordre \(108\) dans le domaine fini déclaré.
\end{theorem}

\begin{proof}
La relation de Weyl \(S^{-1} Z S=\zeta Z\) donne
\(S^{-n} Z S^n=\zeta^n Z\). Sa partie réelle et sa partie imaginaire produisent la
formule. Si un entier \(d>0\) était une période, on aurait
\(\zeta^d=1\) ; la primitivité de \(\zeta\) impose \(N\mid d\).
\end{proof}

\subsection{Stabilité spectrale finie}

La séparation minimale entre racines distinctes de l’unité est
\[
 g_N=\min_{k\ne\ell}|\zeta^k-\zeta^\ell|
 =2\sin\frac{\pi}{N}>0.
\]

\begin{theorem}[Rigidité covariante et stabilité spectrale finie]
\label{thm:hmtf2-estabilidad-cristal}
Les deux affirmations suivantes sont satisfaites.

\begin{enumerate}
 \item Pour toute application unitaire \(V\), le triplet
 \[
  U_F^{(V)}=V U_F V^\dagger,\qquad
  Z^{(V)}=V Z V^\dagger,\qquad
  \rho^{(V)}=V\rho V^\dagger
 \]
 conserve exactement la relation de Weyl, la sous-harmonique d’ordre \(N\) et
 l’amplitude du vecteur d’ordre.
 \item Si \(U'\) est unitaire et
 \(\|U'-U_F\|=\varepsilon<g_N/2\), les \(N\) disques de rayon
 \(g_N/2\) centrés en \(\{\zeta^k\}\) séparent de manière unique les secteurs
 spectraux de \(U'\). Leurs projecteurs de Riesz ont le même rang que
 ceux de \(U_F\) tout au long de toute déformation qui reste dans cette boule.
 En outre, pour tout observable \(A\), tout état \(\rho\) et tout \(n\geq0\),
 \[
 \left|
 \operatorname{Tr}\rho\left(U'^{-n}AU'^n-U_F^{-n}AU_F^n\right)
 \right|
 \leq 2n\varepsilon\|A\|.
 \]
\end{enumerate}
\end{theorem}

\begin{proof}
La première partie s’obtient en conjuguant chaque identité. Pour la seconde, la
séparation \(g_N\) empêche deux secteurs d’entrer dans le même contour ; la
continuité des projecteurs de Riesz maintient leurs rangs. L’identité
télescopique donne
\(\|U'^n-U_F^n\|\leq n\varepsilon\), et la même borne vaut pour les
puissances inverses. L’inégalité des espérances résulte de l’insertion et de la
soustraction de \(U'^{-n}AU_F^n\).
\end{proof}

L’objet obtenu est un \emph{cristal temporel fini HMT} avec forçage,
observable sous-harmonique et stabilité spectrale quantifiée. Qu’un matériau
macroscopique réalise cette classe, conserve la cohérence à grand volume ou
présente une protection préthermique ou localisée relève de la
\textsc{reconnaissance externe} ; ce n’est pas un résidu mathématique de la
construction finie.


